\documentclass[11pt,a4paper]{article}
\usepackage{../../intuitmath}
\hypersetup{pdftitle={The Schr\"odinger Equation as Operator Evolution},pdfauthor={Mathematics and Physics: Intuition First}}
\renewcommand{\sectionmark}[1]{\markboth{#1}{}}
\newcommand{\ket}[1]{\left|#1\right\rangle}
\newcommand{\bra}[1]{\left\langle#1\right|}

\begin{document}
\tutorialtitle{The Schr\"odinger Equation as Operator Evolution}{A bridge from vectors and matrices to wavefunctions}{Why quantum mechanics looks like linear algebra in an infinite-dimensional space}
\tableofcontents
\clearpage

\section{From coordinate lists to functions}

A finite vector may be written
\[
\vect v=(v_1,v_2,\ldots,v_N).
\]
The index $i$ tells us which coordinate we are reading. A function $\psi(x)$ can be viewed similarly: the continuous label $x$ replaces the discrete label $i$.

The finite inner product
\[
\langle\vect u,\vect v\rangle=\sum_i u_i^*v_i
\]
becomes
\[
\boxed{\langle\phi,\psi\rangle
=\int_{-\infty}^{\infty}\phi^*(x)\psi(x)\,\dd x.}
\]
The norm is
\[
\|\psi\|^2=\langle\psi,\psi\rangle
=\int|\psi(x)|^2\,\dd x.
\]

\begin{intuition}[The bridge]
A wavefunction is not merely a graph in physical space. It is one vector in a function space. Its value at $x$ is analogous to a coordinate of that vector in the position-labelled description.
\end{intuition}

\section{Normalisation and probability}

Quantum states are normalised so that
\[
\boxed{\int_{-\infty}^{\infty}|\psi(x,t)|^2\,\dd x=1.}
\]
For an interval $[a,b]$,
\[
\Pr(a\le x\le b)=\int_a^b|\psi(x,t)|^2\,\dd x.
\]
Multiplying $\psi$ by a global phase $\e^{i\alpha}$ does not change $|\psi|^2$, so state vectors that differ only by a global phase represent the same physical ray.

\section{Operators are the infinite-dimensional analogue of matrices}

A matrix maps vectors to vectors. A linear operator maps functions to functions:
\[
\widehat A(\alpha\psi+\beta\phi)
=\alpha\widehat A\psi+\beta\widehat A\phi.
\]
Important examples in one dimension are
\[
\widehat x\psi=x\psi,
\qquad
\widehat p\psi=-i\hbar\pd{\psi}{x}.
\]
The Hamiltonian for a particle in a potential is
\[
\boxed{
\widehat H=-\frac{\hbar^2}{2m}\frac{\partial^2}{\partial x^2}+V(x).}
\]
The kinetic energy is a differential operator; the potential energy is a multiplication operator.

\section{Self-adjointness and real measurements}

An operator is self-adjoint when
\[
\langle\phi,\widehat A\psi\rangle
=\langle\widehat A\phi,\psi\rangle
\]
for all functions in an appropriate domain.

If
\[
\widehat A\psi=a\psi,
\]
then, for a normalised eigenfunction,
\begin{align*}
a&=\langle\psi,a\psi\rangle
=\langle\psi,\widehat A\psi\rangle\\
&=\langle\widehat A\psi,\psi\rangle
=a^*.
\end{align*}
Hence $a$ is real.

\begin{warningbox}[Domains are part of the operator]
A differential expression alone is not a complete operator. Boundary conditions and regularity determine the domain, and the domain determines whether integration by parts produces unwanted boundary terms.
\end{warningbox}

\section{The time-dependent Schr\"odinger equation}

The evolution law is
\[
\boxed{i\hbar\pd{\psi}{t}=\widehat H\psi.}
\]
In vector notation,
\[
i\hbar\td{}{t}\ket{\psi(t)}=\widehat H\ket{\psi(t)}.
\]
This is a first-order differential equation in time for a vector in a function space. Compare it with a finite linear system
\[
\td{\vect v}{t}=A\vect v.
\]
The Hamiltonian is the generator of quantum time evolution.

If $\widehat H$ is time independent, the formal solution is
\[
\boxed{
\ket{\psi(t)}=exp\left(-\frac{i}{\hbar}\widehat Ht\right)\ket{\psi(0)}.}
\]

\section{Why the norm is preserved}

Let $\widehat H$ be self-adjoint. Then
\begin{align*}
\td{}{t}\langle\psi,\psi\rangle
&=\left\langle\pd{\psi}{t},\psi\right\rangle
+\left\langle\psi,\pd{\psi}{t}\right\rangle\\
&=\frac{i}{\hbar}\langle\widehat H\psi,\psi\rangle
-\frac{i}{\hbar}\langle\psi,\widehat H\psi\rangle=0.
\end{align*}
Thus a normalised state stays normalised. The evolution operator is unitary.

\section{Energy eigenfunctions separate the time dependence}

Suppose
\[
\widehat H\phi_n=E_n\phi_n.
\]
Then
\[
\psi_n(x,t)=\phi_n(x)\e^{-iE_nt/\hbar}
\]
solves the time-dependent equation. Its probability density is stationary:
\[
|\psi_n(x,t)|^2=|\phi_n(x)|^2.
\]

For a superposition,
\[
\psi(x,0)=\sum_n c_n\phi_n(x),
\]
linearity gives
\[
\boxed{
\psi(x,t)=\sum_n c_n\e^{-iE_nt/\hbar}\phi_n(x).}
\]
Time evolution rotates each energy coordinate by its own phase.

\section{The infinite square well}

Let
\[
V(x)=0\quad(0<x<L),
\qquad
V=\infty\quad\text{outside}.
\]
The boundary conditions are
\[
\phi(0)=\phi(L)=0.
\]
Inside the well,
\[
-\frac{\hbar^2}{2m}\phi''=E\phi.
\]
Write
\[
k^2=\frac{2mE}{\hbar^2}.
\]
Then
\[
\phi(x)=A\sin(kx)+B\cos(kx).
\]
The condition at $x=0$ gives $B=0$. The condition at $x=L$ gives
\[
\sin(kL)=0
\quad\Longrightarrow\quad
kL=n\pi.
\]
Therefore
\[
\boxed{
\phi_n(x)=\sqrt{\frac2L}\sin\frac{n\pi x}{L},
\qquad
E_n=\frac{\hbar^2\pi^2n^2}{2mL^2}.}
\]

\begin{center}
\begin{tikzpicture}[x=1.0cm,y=.75cm,>=Latex]
  \draw[->,ink!65] (0,0)--(6.5,0) node[right]{$E$};
  \foreach \y/\n in {.65/1,1.55/2,2.75/3,4.25/4}{
    \draw[blue,line width=1.6pt] (1.2,\y)--(4.7,\y);
    \node[left,ink] at (1.2,\y){$E_{\n}$};
  }
  \node[ink,align=left] at (5.4,2.4){spacing grows\\because $E_n\propto n^2$};
\end{tikzpicture}
\end{center}

Quantisation is a boundary-condition effect: only certain standing-wave shapes fit the interval.

\section{Fourier coordinates diagonalise free evolution}

For a free particle,
\[
\widehat H=-\frac{\hbar^2}{2m}\partial_x^2.
\]
The plane wave $\e^{ikx}$ is an eigenfunction because
\[
\partial_x^2\e^{ikx}=-k^2\e^{ikx}.
\]
Thus
\[
\widehat H\e^{ikx}=\frac{\hbar^2k^2}{2m}\e^{ikx}.
\]
In Fourier coordinates, the differential operator becomes multiplication by $\hbar^2k^2/(2m)$. Each Fourier component evolves independently:
\[
\widetilde\psi(k,t)
=\e^{-i\hbar k^2t/(2m)}\widetilde\psi(k,0).
\]

\section{Heat flow and quantum evolution}

The heat equation
\[
\pd{u}{t}=D\partial_x^2u
\]
and the free Schr\"odinger equation
\[
\pd{\psi}{t}=\frac{i\hbar}{2m}\partial_x^2\psi
\]
use the same spatial operator but different time generators.

For a Fourier mode $\e^{ikx}$:
\[
\widetilde u(k,t)=\e^{-Dk^2t}\widetilde u(k,0),
\]
so high-frequency heat modes decay. Quantum modes acquire phases:
\[
\widetilde\psi(k,t)=\e^{-i\hbar k^2t/(2m)}\widetilde\psi(k,0),
\]
so their magnitudes do not decay.

\begin{resultbox}[Same basis, different geometry]
Heat flow contracts vectors toward smoother states. Quantum evolution rotates state-space coordinates without changing the norm.
\end{resultbox}

\section{Commutators and incompatible operations}

The commutator is
\[
[\widehat A,\widehat B]=\widehat A\widehat B-\widehat B\widehat A.
\]
For position and momentum,
\begin{align*}
[\widehat x,\widehat p]\psi
&=x(-i\hbar\psi')+i\hbar(x\psi)'\\
&=i\hbar\psi.
\end{align*}
Hence
\[
\boxed{[\widehat x,\widehat p]=i\hbar\widehat I.}
\]
Noncommutation is the operator statement behind the position--momentum uncertainty relation.

\section{Expectation values}

For a normalised state,
\[
\boxed{\langle A\rangle
=\langle\psi,\widehat A\psi\rangle
=\int\psi^*(x)(\widehat A\psi)(x)\,\dd x.}
\]
The spread is
\[
(\Delta A)^2=\langle(\widehat A-\langle A\rangle)^2\rangle.
\]
The expectation is a weighted average over possible measurement outcomes, not generally an eigenvalue attached to the current state.

\section{What to remember}

\begin{enumerate}
\item A wavefunction is a vector in a function space.
\item Integrals replace finite coordinate sums in inner products.
\item Differential operators play the role of matrices.
\item Self-adjointness supplies real eigenvalues and unitary evolution.
\item The Hamiltonian generates time evolution.
\item Choosing an eigenbasis can turn an operator equation into independent scalar equations for coordinates.
\end{enumerate}

\begin{checkBox}[The infinite-dimensional bridge]
The conceptual jump from linear algebra to quantum mechanics is smaller than it first appears: replace finite coordinate lists by functions, sums by integrals, matrices by operators, and matrix exponentials by operator exponentials.
\end{checkBox}

\end{document}
